\section{Development priorities and versioned guarantees}\label{sec:evolution}

Four identifiers remain explicit: manuscript version, software/source identity, certificate schema, and semantic contract. Manuscript 3.0 describes the technical snapshot through PR51 and accompanies software \texttt{0.3.0a1} research-alpha preparation. Coordinatewise schemas remain unchanged; research profiles retain individual schemas and source bindings. A release tag or deposit identifier must refer to the actual published revision, while a concept DOI denotes an evolving record. The existing version-2 Paper III DOI must not be presented as a new version-3 identifier.

Several tasks that were future work in version 2 now have concrete implementations: restricted automatic observation inference, source-bound signed execution, joint-carrier caller checks, initial pure-row semantics, source-to-consumer ground proofs, checked lemmas, and applicable summaries. F1a/F1b add general formal soundness for decoded positive proofs and chronological composition. Each closes a specific obligation under its stated contract; none establishes an unrestricted source-language verifier.

The next implementation choice should preserve the useful simpler ordinary-SDK path while testing additional mechanisms against it under equal delivery contracts. The registered broader comparison remains future work. It must retain native policies, caches, complete first-proof costs, cold receiver checks, application/assessment distinctions, failures, and memory tradeoffs. A smaller factor, fewer questions, fewer calls, or fewer bytes is informative but does not alone establish lower total cost.

Observation inference remains incomplete beyond the current finite language and restricted residual models. The next expansion should be justified by a concrete consumer that current native rules cannot handle, with explicit source correspondence and a frozen external evaluation after the new engine is fixed. Known development successes should not be relabeled as unseen external coverage. The finite-ground query kernel already distinguishes proof visibility from insufficient premises; a new candidate language must preserve that distinction rather than interpreting every negative model as a source bug.

The next formal obligations are the native source rules actually used by successful consumers, the source-to-term adapter, and the interface-to-runtime correspondence. For the Graal line, the descending and concrete caller bridges remain separate from the conditional mathematical composition theorems. F1 does not resolve these by itself. Coverage, discovery effort, portable proof size, replay cost, runtime cost, native correspondence, and formalized semantic scope remain separate progress measures.
